% =============================================================================
% Section 10: Critical Path
% =============================================================================

\section{Critical Path}
\sectionauthors{CM, JB}

\subsection{Critical Path Definition}

The critical path is the longest sequence of dependent activities that determines the minimum project completion time. Activities on the critical path have zero slack (float), meaning any delay in a critical path activity directly delays the project end date.

For the STAR Robot project, the critical path spans 214 working days and includes 47 activities across all project phases.

\subsection{Critical Path Summary}

\begin{table}[H]
\centering
\caption{Critical Path Statistics}
\label{tab:critical_path_stats}
\begin{tabular}{@{}lr@{}}
\toprule
\textbf{Metric} & \textbf{Value} \\
\midrule
Total Activities on Critical Path & 47 \\
Critical Path Duration & 214 days \\
Total Project Activities & 84 \\
Percentage on Critical Path & 56\% \\
Total Float (Non-Critical Activities) & Varies (1-45 days) \\
\bottomrule
\end{tabular}
\end{table}

\subsection{Critical Path Visualization}

The critical path is highlighted in red on the NLD (see Figure \ref{fig:nld_critical_path} in Section \ref{sec:nld_sec}) and the Gantt chart. Full diagrams are available in the Appendix.

\subsection{Critical Path Activities}

The following activities comprise the critical path, organized by phase. Row colors match the phase color coding used in the NLD and WBS diagrams.

\subsubsection{Research Phase Critical Activities}

\begin{table}[H]
\centering
\caption{Research Phase Critical Activities}
\label{tab:critical_path_research}
\begin{tabular}{@{}llr@{}}
\toprule
\textbf{WBS ID} & \textbf{Activity} & \textbf{Duration} \\
\midrule
\rowcolor{researchblue!30} 1.1.1.1 & Gather stakeholder requirements & 5 days \\
\rowcolor{researchblue!30} 1.1.1.2 & Document functional requirements & 5 days \\
\rowcolor{researchblue!30} 1.1.2.1 & Research SLAM algorithms & 4 days \\
\rowcolor{researchblue!30} 1.1.2.2 & Evaluate sensor technologies & 4 days \\
\rowcolor{researchblue!30} 1.1.3.1 & Select compute platform & 5 days \\
\rowcolor{researchblue!30} 1.1.3.2 & Select sensors & 4 days \\
\rowcolor{researchblue!30} 1.1.6.1 & Compile research findings & 5 days \\
\rowcolor{researchblue!30} 1.1.6.2 & Create research report & 5 days \\
\bottomrule
\end{tabular}
\end{table}

\subsubsection{Design Phase Critical Activities}

\begin{table}[H]
\centering
\caption{Design Phase Critical Activities}
\label{tab:critical_path_design}
\begin{tabular}{@{}llr@{}}
\toprule
\textbf{WBS ID} & \textbf{Activity} & \textbf{Duration} \\
\midrule
\rowcolor{designyellow!30} 1.2.1.1 & Define system architecture & 5 days \\
\rowcolor{designyellow!30} 1.2.1.2 & Create interface specifications & 5 days \\
\rowcolor{designyellow!30} 1.2.1.3 & Document design decisions & 5 days \\
\rowcolor{designyellow!30} 1.2.2.1 & Design chassis frame & 5 days \\
\rowcolor{designyellow!30} 1.2.2.2 & Design drive system & 5 days \\
\rowcolor{designyellow!30} 1.2.2.3 & Design sensor mounts & 5 days \\
\rowcolor{designyellow!30} 1.2.2.4 & Create CAD assembly & 5 days \\
\rowcolor{designyellow!30} 1.2.6.1 & Prepare design review & 5 days \\
\rowcolor{designyellow!30} 1.2.6.2 & Conduct design review & 5 days \\
\bottomrule
\end{tabular}
\end{table}

\subsubsection{Build Phase Critical Activities}

\begin{table}[H]
\centering
\caption{Build Phase Critical Activities}
\label{tab:critical_path_build}
\begin{tabular}{@{}llr@{}}
\toprule
\textbf{WBS ID} & \textbf{Activity} & \textbf{Duration} \\
\midrule
\rowcolor{buildorange!30} 1.3.1.1 & Fabricate chassis components & 5 days \\
\rowcolor{buildorange!30} 1.3.1.2 & Assemble chassis frame & 5 days \\
\rowcolor{buildorange!30} 1.3.1.3 & Install drive system & 5 days \\
\rowcolor{buildorange!30} 1.3.1.4 & Verify mechanical assembly & 5 days \\
\rowcolor{buildorange!30} 1.3.3.1 & Install LiDAR sensor & 5 days \\
\rowcolor{buildorange!30} 1.3.3.2 & Install depth camera & 5 days \\
\rowcolor{buildorange!30} 1.3.3.3 & Install compute platform & 5 days \\
\rowcolor{buildorange!30} 1.3.3.4 & Configure sensor interfaces & 5 days \\
\rowcolor{buildorange!30} 1.3.3.5 & Verify sensor operation & 5 days \\
\rowcolor{buildorange!30} 1.3.4.1 & Develop motor control firmware & 5 days \\
\rowcolor{buildorange!30} 1.3.4.2 & Develop ROS 2 interface & 5 days \\
\rowcolor{buildorange!30} 1.3.5.1 & Integrate hardware subsystems & 5 days \\
\rowcolor{buildorange!30} 1.3.5.2 & Integrate software components & 5 days \\
\rowcolor{buildorange!30} 1.3.5.3 & System bring-up testing & 5 days \\
\rowcolor{buildorange!30} 1.3.6.1 & Verify build completeness & 5 days \\
\rowcolor{buildorange!30} 1.3.6.2 & Document build status & 5 days \\
\bottomrule
\end{tabular}
\end{table}

\subsubsection{Test Phase Critical Activities}

\begin{table}[H]
\centering
\caption{Test Phase Critical Activities}
\label{tab:critical_path_test}
\begin{tabular}{@{}llr@{}}
\toprule
\textbf{WBS ID} & \textbf{Activity} & \textbf{Duration} \\
\midrule
\rowcolor{testgreen!30} 1.4.1.1 & Test sensor subsystems & 5 days \\
\rowcolor{testgreen!30} 1.4.1.2 & Test motor control & 5 days \\
\rowcolor{testgreen!30} 1.4.1.3 & Test software modules & 5 days \\
\rowcolor{testgreen!30} 1.4.2.1 & Test SLAM performance & 5 days \\
\rowcolor{testgreen!30} 1.4.2.2 & Test navigation functions & 5 days \\
\rowcolor{testgreen!30} 1.4.2.3 & Test obstacle avoidance & 5 days \\
\rowcolor{testgreen!30} 1.4.2.4 & Document integration results & 5 days \\
\rowcolor{testgreen!30} 1.4.4.1 & Measure system accuracy & 5 days \\
\rowcolor{testgreen!30} 1.4.4.2 & Measure system latency & 5 days \\
\rowcolor{testgreen!30} 1.4.5.1 & Conduct acceptance tests & 5 days \\
\rowcolor{testgreen!30} 1.4.5.2 & Document test results & 5 days \\
\bottomrule
\end{tabular}
\end{table}

\subsubsection{Close Phase Critical Activities}

\begin{table}[H]
\centering
\caption{Close Phase Critical Activities}
\label{tab:critical_path_close}
\begin{tabular}{@{}llr@{}}
\toprule
\textbf{WBS ID} & \textbf{Activity} & \textbf{Duration} \\
\midrule
\rowcolor{closepink!30} 1.5.1.1 & Complete user documentation & 5 days \\
\rowcolor{closepink!30} 1.5.1.2 & Complete technical documentation & 5 days \\
\rowcolor{closepink!30} 1.5.3 & Deliver final report & 5 days \\
\bottomrule
\end{tabular}
\end{table}

\subsection{Critical Path Analysis}

\subsubsection{Why These Activities Are Critical}

The critical path activities share these characteristics:

\begin{itemize}
    \item \textbf{Sequential Dependencies}: Each activity depends on completion of the previous activity
    \item \textbf{No Parallel Alternatives}: Work cannot be redistributed to other paths
    \item \textbf{Resource Constraints}: Key team members are required for specific expertise
    \item \textbf{Technical Dependencies}: Outputs are required inputs for subsequent work
\end{itemize}

\subsubsection{Critical Path Risk Factors}

The following factors pose the greatest risk to critical path activities:

\begin{enumerate}
    \item \textbf{Component Procurement}: Delays in receiving LiDAR or depth camera affect sensor integration
    \item \textbf{SLAM Performance}: Algorithm tuning may require iteration if accuracy targets not met
    \item \textbf{Integration Complexity}: Hardware-software integration often reveals unexpected issues
    \item \textbf{Testing Failures}: Failed tests require debugging and re-testing cycles
\end{enumerate}

\subsection{Critical Path Management}

\subsubsection{Monitoring}

Critical path status is monitored through:

\begin{itemize}
    \item Weekly progress review against baseline schedule
    \item Early warning indicators from in-progress activities
    \item Dependency tracking for predecessor completion
\end{itemize}

\subsubsection{Mitigation Strategies}

To protect the critical path:

\begin{itemize}
    \item \textbf{Buffer Time}: Non-critical activities provide schedule buffer if resources shift
    \item \textbf{Parallel Development}: Software development proceeds parallel to hardware
    \item \textbf{Early Procurement}: Long-lead items ordered during Research phase
    \item \textbf{Resource Flexibility}: Cross-trained team members can support critical activities
    \item \textbf{Risk Reserves}: Schedule contingency available for high-risk activities
\end{itemize}

\subsubsection{Recovery Options}

If critical path delays occur:

\begin{enumerate}
    \item \textbf{Crashing}: Add resources to critical activities (limited effectiveness)
    \item \textbf{Fast-Tracking}: Overlap sequential activities where possible
    \item \textbf{Scope Reduction}: Defer non-essential features to recover schedule
    \item \textbf{Schedule Compression}: Reduce activity durations through focused effort
\end{enumerate}

\subsection{Float Analysis}

Activities not on the critical path have float (slack) time:

\begin{table}[H]
\centering
\caption{Non-Critical Activities with Highest Float}
\label{tab:float_analysis}
\begin{tabular}{@{}llr@{}}
\toprule
\textbf{Activity} & \textbf{Phase} & \textbf{Float (Days)} \\
\midrule
\rowcolor{researchblue!30} Risk Assessment & Research & 15 \\
\rowcolor{designyellow!30} Software Architecture & Design & 12 \\
\rowcolor{designyellow!30} Electrical Design & Design & 10 \\
\rowcolor{gray!30} PM Documentation & PM Overhead & 45 \\
\rowcolor{closepink!30} Knowledge Transfer & Close & 8 \\
\bottomrule
\end{tabular}
\end{table}

Float provides schedule flexibility and opportunities to redirect resources to critical path activities if needed.
